\documentclass[11pt]{amsart}
\usepackage{amsmath,amssymb,amsthm}
\usepackage{booktabs}
\usepackage[margin=1in]{geometry}
\usepackage{hyperref}
\hypersetup{pdftitle={The Halting Problem of cm16/26/51/52},pdfauthor={CyclingCycloid}}

\newtheorem{theorem}{Theorem}[section]
\newtheorem{lemma}[theorem]{Lemma}
\newtheorem{definition}[theorem]{Definition}
\newtheorem{remark}[theorem]{Remark}
\newtheorem{conjecture}[theorem]{Conjecture}

\newcommand{\HALT}{\textsf{HALT}}

\title{The Halting Problem of the Four Turing Machines \texttt{cm16/26/51/52}:\\ Criteria, Era Algebra, Correspondence Facts, and the Open Problem}
\author{}
\date{October 2026}

\begin{document}
\maketitle

\section{The four machines and conventions}

\begin{center}\small
\begin{tabular}{lll}
\toprule
Machine & Transition table & Halting transition\\
\midrule
cm16 & \texttt{1RB0LE\_0RC1RB\_1LC1RD\_1RF0LA\_0RD1LE\_1RB---} & $F$ on $0$\\
cm26 & \texttt{1RB1RC\_1LC0LE\_1RD1RC\_1LB0RA\_---0LF\_1RD1LF} & $E$ on $0$\\
cm51 & \texttt{1RB1RA\_1LC0RC\_1RE0LD\_0RE0LF\_1RA---\_1RB1LF} & $E$ on $0$\\
cm52 & \texttt{1RB0LD\_1LC1RD\_0LE1LA\_0RB1RD\_---0LF\_1RA1LF} & $D$ on $0$\\
\bottomrule
\end{tabular}
\end{center}

Bi-infinite tape, initially all $0$, head at cell $0$, state $A$. The four halting transitions occupy different positions; criteria are established independently per machine (\S2), followed by the common era algebra (\S3), the verified correspondence facts (\S4), and the per-machine reduction and open problem (\S5--\S6). All criterion- and routine-layer assertions are exhaustively verified by two checkers: \texttt{casecheck-halt.mjs} (204,800 configurations, 111,257 halting instances, zero counterexamples, entry enumeration ENUM-OK $\times$4) and \texttt{casecheck-island.mjs} (39 assertions, all passing).

\section{Halt criteria (independent per machine)}

\begin{theorem}[cm16]\label{thm:h16}
cm16 halts if and only if at some time state $D$ reads a $0$ at cell $w$ with $\tau(w{+}1)=1$.
\end{theorem}
\begin{theorem}[cm26]\label{thm:h26}
cm26 halts if and only if at some time state $B$ reads a $1$ at cell $q$ with $\tau(q{-}1)=0$.
\end{theorem}
\begin{theorem}[cm51]\label{thm:h51}
cm51 halts if and only if at some time state $C$ or $D$ reads a $0$ at cell $w$ with $\tau(w{+}1)=1$.
\end{theorem}
\begin{theorem}[cm52]\label{thm:h52}
cm52 halts if and only if at some time state $C$ reads a $0$ at cell $c$ with $\tau(c{-}1)=0$.
\end{theorem}

\begin{proof}[Method (identical form for all four)]
Halting $\iff$ reaching the machine's halting state reading $0$. Enumerate at transition level all in-edges of the halting state (as in Theorem 4.1 of the cm2 paper: trace the provenance of the cell values read), reducing ``halting state reads $0$'' to the visit of a specific two-cell local pattern; then instantiate both directions over all configurations in the pattern window: configurations satisfying the pattern halt within finitely many steps, the others leave the window along their routine branches. Exhaustiveness via the checker: frozen skeleton (per machine, from deep scans) $\times$ working-triple boundaries $\times$ head classes, 204,800 cells, both directions, zero counterexamples [casecheck-halt].
\end{proof}

\section{Lap grammar and era algebra}

\begin{definition}[marker states and eras]
A marker state is the $(A,h,\tau)$ at the end of a dance-family routine; the evolution between consecutive markers is one era. The blank-tape initial segment enters the first marker by direct simulation.
\end{definition}

\begin{lemma}[standard lap word]\label{lem:lap}
For each machine and each parameter $R\ge1$, the transition word of one lap macro step is
\[
A0\,F0 \;\Big|\; B1^{R-2}\,B0\,C0^2\,C1\,D1\,A1 \;\Big|\; E1^{R-1}\,E0\,D1\,A0,
\]
$2R+14$ steps in total, with net effect: right edge of the left reservoir $-1$, right reservoir $-2$, a deposit three cells away $+1$ (bookkeeping $(x{+}1,\,R{-}2,\,z{+}1)$), head one cell left.
\end{lemma}
\begin{proof}
The word is read off the transition table step by step; the parameter branches (residue classes of $R$ and boundaries $R\in\{1,2,3,4\}$) are verified instance-by-instance for $R\in\{1,\dots,40\}$ and against each machine's deep-scan lap anchors (about 780,000 laps in total) by cell-diff [casecheck-island].
\end{proof}

\begin{definition}[eras and resets]
One era = the lap chain up to reservoir exhaustion + the drain routine + the reset. Reservoir parameters $(L,M,R)$ evolve along the chain by Lemma~\ref{lem:lap}; draining triggers a churn cascade (reseeding). Era key sets (key = (position class, neighborhood block count)) collected from deep scans: cm16/26/51 share 78 keys identically; cm52 has 89.
\end{definition}

\section{The era system: data and structure}

\begin{definition}[reservoir configuration]
The marker tape consists of: a left reservoir $1$-block, a right reservoir $1$-block and a deposit block, separated by $0^2$ seams, with the frozen skeleton as tail (identical bit-for-bit across the four machines, fixed by deep scans); the working triple $(x,R,z)$ records the working-block, right-reservoir and deposit lengths.
\end{definition}

\begin{remark}[era-chain example (cm16, near $10^9$ steps)]
The measured working-triple sequence of eras:
\[
(5392,\,22850,\,1818)\ \longrightarrow\ (1787,\,16815,\,1)\ \xrightarrow{\ \text{churn}\ }\ (14188,\,11225,\,946)\ \longrightarrow\ \cdots
\]
The right reservoir decreases by steps of 2 along the lap chain until exhaustion, triggering the drain and the churn cascade (reseeding into a new reservoir layout); the chain truncations and reset branches form the era key set.
\end{remark}

\begin{remark}[cross-machine alignment data]
The large-era $(L_0,R_0)$ sequences of the four machines align over 562 measured pairs at constant offsets ($0$ or $+1$) with zero misalignment (deep-scan prefix, at the $5\times10^9$-step scale); the era key sets are 78 keys, identical key-by-key, for cm16/26/51 and 89 for cm52 (due to criterion and skeleton differences). These are observations (grade B), not used as theorems.
\end{remark}

\section{Verified correspondence facts between the four machines}

\begin{remark}
The following are layered facts, not an isomorphism theorem (the four transition tables differ; no transition-level isomorphism exists):
\begin{enumerate}
\item the lap grammar and bookkeeping are the same on all four machines (Lemma~\ref{lem:lap}, verified per machine);
\item the large-era $(L_0,R_0)$ sequences align over 562 pairs at constant offsets ($0$ or $+1$) with zero misalignment --- an observation on the deep-scan prefix [at the $5\times10^9$-step scale];
\item end configurations agree bit-for-bit only on the frozen skeleton; working triples differ by $\pm100$ (a correction of the earlier ``bitwise identical'' measurement);
\item the halt criteria have different sources (Theorems~\ref{thm:h16}--\ref{thm:h52}); non-halting assembly proceeds independently per machine.
\end{enumerate}
\end{remark}

\section{The island invariant and the danger pattern}

\begin{definition}[island]
An island is the pattern $(0^1,1^1,0^1)$ (a single-cell $1$-island). A danger event is an occurrence of this pattern with the \textbf{head far} (not co-located with the destroyer wave).
\end{definition}

\begin{theorem}[island invariant]\label{thm:island}
Within the verified domain ($2\times10^7$ steps per machine, 892 island moments): every island occurrence is co-located with the destroyer wave (head distance $\le1$); island occurrences are transient --- within each lap islands exist in \emph{two} segments (correcting the earlier ``only the boundary $R{=}1$ segment''), with creator and destroyer the same passing wave and existence window shorter than one lap [casecheck-island, 39 assertions].
\end{theorem}

\begin{definition}[vocabulary, audit and guard]
The \emph{verified vocabulary} is the set of routine branches actually occurring and catalogued in deep scans and censuses (the standard lap word with its parameter branches, drain branches, churn reseed branches; the authoritative inventory is the era data and the classifier of \texttt{quad.c}). \emph{Audit mode}: at every $A$-state landing, check whether the danger pattern (a head-far $(0^1,1^1,0^1)$) appears in the five-cell window (head neighborhood $\pm2$), stepping to a given depth. \emph{Guard mode}: at every step, check directly whether the machine's criterion (the two-cell condition of Theorems~\ref{thm:h16}--\ref{thm:h52}) holds, stepping to a given depth.
\end{definition}

\begin{theorem}[safety within the vocabulary]\label{thm:vocab}
At every $A$-state landing covered by the verified routine vocabulary, the pattern preconditions of Theorems~\ref{thm:h16}--\ref{thm:h52} do not occur (sprint-landing right gaps $\ge2$; island positions excluded by Theorem~\ref{thm:island}).
\end{theorem}

\section{Reduction and open problem}

\begin{theorem}[reduction, per machine]\label{thm:reduction}
For each machine $m\in\{16,26,51,52\}$ the following are equivalent:
\begin{enumerate}
\item $m$ halts;
\item the criterion pattern of $m$ (the two-cell condition of Theorems~\ref{thm:h16}--\ref{thm:h52}) holds at an \emph{out-of-vocabulary} landing (a danger event).
\end{enumerate}
Machine verification: audit mode steps through five-cell checks at out-of-vocabulary landings ($10^9$ steps $\times$4) with zero danger occurrences; guard mode ($2\times10^{11}$ steps $\times$4, stepwise halt-precondition checking) with zero triggers; bare-machine deep scan of cm16 to $5\times10^{11}$ steps with zero halts.
\end{theorem}
\begin{proof}
(2)$\Rightarrow$(1): the criterion pattern triggers the in-edge chain of Theorems~\ref{thm:h16}--\ref{thm:h52}; the machine halts. (1)$\Rightarrow$(2): suppose a halt; then some criterion pattern is visited; by Theorem~\ref{thm:vocab} (vocabulary safety) the visit occurs at an $A$-state landing outside the vocabulary coverage; the enumeration completeness of out-of-vocabulary landings is carried by the census/era key coverage [A$^-$].
\end{proof}

\begin{conjecture}[five-cell avoidance]\label{conj:quad}
No danger event ever occurs in the runs of the four machines (formally: clause (2) of Theorem~\ref{thm:reduction} is false for all time).
\end{conjecture}

Refutation record: the local-abstraction route (key graphs) is dead --- in the radius-ladder experiment for key-set successor closure, $|K^*|$ grows by about $\times3.1$ per added radius (cm16: $99\to294\to939\to2863\to8879$ for $r=2..6$; two reruns reproduce digit-for-digit), with danger-key hits growing in step, so no finite fixed point exists for fixed-radius local abstraction [kclosure]. Research field: piecewise-affine trajectory avoidance of the same shape as the cm2/M1 residuals (classification by analogy; per-machine formal restatement not performed).

\section*{Reproduction}
\texttt{node casecheck-halt.mjs \&\& node casecheck-island.mjs \&\& node kclosure.mjs}; \texttt{quad6.exe <code> audit 1000000000}; \texttt{quad6.exe <code> guard 200000000000}. All one-line commands.

\end{document}
